% Çıkarım koordinatları fiziğe NEREDEN giriyor?
% (a) Lundborg akma yüzeyi: Y0 yalnız DÜŞÜK basınçta ayırt edilebiliyor.
% (b) P-alpha ezilme eğrisi: alpha_b bloğun ne kadar sıkışabileceğini
%     belirliyor.
% Eğriler elle çizilmedi; veri kodun kendi formülünden üretildi
% (scripts/rapor2204_veri.py -> veri/lundborg.dat, veri/crush.dat).
% `docs/rapor2204` içinden derle:  lualatex sekil/malzeme.tex
\documentclass[border=4pt,tikz]{standalone}
\input{sekil/ortak.tex}
\usepackage[outline]{contour}
\contourlength{1.1pt}
\tikzset{>=latex}

\colorlet{ink}{black!82}
\colorlet{dusuk}{blue!55!cyan!78!black}   % Y0 = 10 Pa
\colorlet{orta}{violet!60!black}          % Y0 = 500 Pa
\colorlet{yuksek}{vurgu}                  % Y0 = 5000 Pa
\colorlet{matriscol}{koyu!55}

\pgfplotstableread{veri/lundborg.dat}\lund
\pgfplotstableread{veri/crush.dat}\crush

\begin{document}
\begin{tikzpicture}[font=\footnotesize]

% ===================== (a) akma yüzeyi =================================
\begin{axis}[
  name=panA,
  width=66mm, height=50mm,
  xmode=log, ymode=log,
  xmin=1e1, xmax=1e11, ymin=5, ymax=2e10,
  xlabel={basınç $P$\:[Pa]},
  ylabel={akma dayanımı $Y$\:[Pa]},
  xtick={1e2,1e5,1e8,1e11},
  xticklabels={$10^{2}$,$10^{5}$,$10^{8}$,$10^{11}$},
  ytick={1e1,1e3,1e5,1e7,1e9},
  yticklabels={$10$,$10^{3}$,$10^{5}$,$10^{7}$,$10^{9}$},
  axis line style={draw=ink!65}, tick style={draw=ink!65},
  ticklabel style={scale=0.86, text=ink!80},
  label style={scale=0.9, text=ink},
  title={(a) akma yüzeyi: $Y_0$ nereye giriyor},
  title style={scale=0.95, text=ink, yshift=2pt},
  axis on top=false,
]
  % Üç eğrinin pratikte ayrışamaz hâle geldiği bölge.
  \fill[ink!7] (axis cs:1e5,5) rectangle (axis cs:1e11,2e10);
  \node[scale=0.76, text=ink!60, anchor=south east, align=right, inner sep=3pt]
       at (axis cs:0.9e11,8)
       {\contour{ink!7}{üç eğri ayrışamaz}\\[-1pt]
        \contour{ink!7}{(fark $<\%1$)}};

  % YM doygunluk tavanı
  \draw[densely dashed, ink!45, line width=0.45pt]
       (axis cs:1e1,1.5e9) -- (axis cs:1e11,1.5e9);
  % Sağa alındı: sol üst köşede eksen çerçevesine giriyordu.
  \node[scale=0.76, text=ink!65, anchor=south west, inner sep=2pt]
       at (axis cs:2e1,1.75e9) {doygunluk $Y_M=1{,}5\times10^{9}$ Pa};

  \addplot[yuksek, line width=0.9pt] table[x=P, y=Y5000] {\lund};
  \addplot[orta,   line width=0.9pt] table[x=P, y=Y500]  {\lund};
  \addplot[dusuk,  line width=0.9pt] table[x=P, y=Y10]   {\lund};

  % P -> 0 sınırında Y = Y0: çıkarımın tuttuğu yer tam burası.
  % Üçü elle yazıldı: \foreach döngü değişkeni renk seçeneği olarak
  % geçirilince pgfkeys onu anahtar adı sanıyor. Ayrıca \c döngü değişkeni
  % olarak KULLANILAMAZ — LaTeX'in sedilla komutudur, Türkçe "ç" patlar.
  \fill[yuksek] (axis cs:1e1,5000) circle (1.4pt);
  \fill[orta]   (axis cs:1e1,500)  circle (1.4pt);
  \fill[dusuk]  (axis cs:1e1,10)   circle (1.4pt);
  % Etiketler eğrilerin üstüne düştüğü için beyaz konturlu.
  \node[scale=0.8, text=yuksek, anchor=west, inner sep=3pt]
       at (axis cs:1.5e1,6400) {\contour{white}{$Y_0=5000$ Pa}};
  \node[scale=0.8, text=orta, anchor=west, inner sep=3pt]
       at (axis cs:1.5e1,520) {\contour{white}{$500$}};
  \node[scale=0.8, text=dusuk, anchor=west, inner sep=3pt]
       at (axis cs:1.5e1,9) {\contour{white}{$10$}};

  % Eğim etiketi. Açı ölçekten: 10 ondalık 66 mm'ye, 9,6 ondalık 50 mm'ye
  % düşüyor; log-log'da 1 eğimli doğru atan(5,2/6,6) = 38 derece görünür.
  \node[scale=0.78, text=ink!70, anchor=south west, inner sep=2pt,
        rotate=38] at (axis cs:1.2e5,1.1e5)
       {\contour{white}{eğim $\mu_f=0{,}6$}};
\end{axis}

% ===================== (b) ezilme eğrisi ===============================
\begin{axis}[
  name=panB,
  at={($(panA.east)+(21mm,0)$)}, anchor=west,
  width=66mm, height=50mm,
  xmode=log,
  xmin=1e4, xmax=1e9, ymin=0.97, ymax=1.62,
  xlabel={basınç $P$\:[Pa]},
  ylabel={distansiyon $\alpha=\rho_s/\rho$},
  xtick={1e4,1e5,1e6,1e7,1e8,1e9},
  xticklabels={$10^{4}$,,$10^{6}$,,$10^{8}$,},
  ytick={1.0,1.1,1.2,1.3,1.4,1.5,1.6},
  yticklabels={$1{,}0$,$1{,}1$,$1{,}2$,$1{,}3$,$1{,}4$,$1{,}5$,},
  axis line style={draw=ink!65}, tick style={draw=ink!65},
  ticklabel style={scale=0.86, text=ink!80},
  label style={scale=0.9, text=ink},
  title={(b) ezilme eğrisi: $\alpha_b$ nereye giriyor},
  title style={scale=0.95, text=ink, yshift=2pt},
  axis on top=false,
]
  % Pe ve Ps eşikleri
  \draw[densely dashed, ink!45, line width=0.45pt]
       (axis cs:1e6,0.97) -- (axis cs:1e6,1.62);
  \draw[densely dashed, ink!45, line width=0.45pt]
       (axis cs:1e8,0.97) -- (axis cs:1e8,1.62);
  \node[scale=0.76, text=ink!65, anchor=south, inner sep=2pt]
       at (axis cs:1e6,1.555) {\contour{white}{$P_e$}};
  \node[scale=0.76, text=ink!65, anchor=south, inner sep=2pt]
       at (axis cs:1e8,1.555) {\contour{white}{$P_s$}};

  \addplot[matriscol, line width=0.9pt, densely dashdotted]
    table[x=P, y=matris] {\crush};
  \addplot[yuksek, line width=0.9pt] table[x=P, y=ab130] {\crush};
  \addplot[orta,   line width=0.9pt] table[x=P, y=ab115] {\crush};
  \addplot[dusuk,  line width=0.9pt] table[x=P, y=ab100] {\crush};

  \node[scale=0.8, text=matriscol, anchor=west, inner sep=3pt]
       at (axis cs:1.3e4,1.52) {matris, $\alpha_0=1{,}5$};
  \node[scale=0.8, text=yuksek, anchor=west, inner sep=3pt]
       at (axis cs:1.3e4,1.325) {$\alpha_b=1{,}30$};
  \node[scale=0.8, text=orta, anchor=west, inner sep=3pt]
       at (axis cs:1.3e4,1.175) {$1{,}15$};
  \node[scale=0.8, text=dusuk, anchor=west, inner sep=3pt]
       at (axis cs:1.3e4,1.025) {$1{,}00$ (gözeneksiz)};

  \node[scale=0.76, text=ink!65, anchor=north east, align=right, inner sep=3pt]
       at (axis cs:0.9e9,1.60)
       {\contour{white}{$P_s$ üstünde boşluk}\\[-1pt]
        \contour{white}{tümüyle kapanır}};
\end{axis}

\end{tikzpicture}
\end{document}
